% TARI29_Assignment_report_template.tex
% This is a template for the assignment reports
% TARI29 Artificial Intelligence, 7.5 credits, Autumn 2025
% Original version by:  Vladimir Tarasov, 2024
% Revised by:           Alexandros Tzanetos, 2024 

% Use a modern class instead of an old article class
\documentclass{scrartcl}

\KOMAoptions{
    parskip=half,  % full, off
    fontsize=12pt, % base font size (10pt default)
    % headings=big,% small/normal/big headings (normal is default), 
    % paper=a5,    % paper format (a4 default) 
    pagesize=auto  % Use paper format for PDF too
}

% ---------------------- Report details --------------------- %
\newcommand{\docauthor}{Andrej Kutny, Charles Madjeri, Samrat Debnath, \\ Vladimir Giustacchini, Aikeya Ainiwaer}
\newcommand{\courseyear}{2025}
\newcommand{\coursename}{TARI29 Artificial Intelligence}
\newcommand{\coursecredits}{7.5 credits}
\newcommand{\fullcoursename}{\coursename, \coursecredits}
\newcommand{\termname}{Autumn \courseyear}
\newcommand{\groupnumber}{4}
\newcommand{\reportname}{Group \groupnumber\ Assignment Report}
\newcommand{\assignmentname}{Evolutionary Computation Assignment}
% ------------------ End of report details ------------------ %



% ---------- Setting up properties of the PDF file ---------- %
\usepackage{hyperref}
\hypersetup{
    colorlinks=true,
    allcolors=blue,
    pdftitle={\reportname},
    pdfauthor={\docauthor},
    pdfsubject={\coursename, \termname}
	% pdfpagemode=FullScreen,
	% pdfpagemode= UseOutlines % the default if present
}
% ------------ End of properties of the PDF file ------------ %



% ----------- Setting up the headers and footers ------------ %
\usepackage{scrlayer-scrpage}
\pagestyle{scrheadings}
\KOMAoptions{% 
    headsepline,% line below the header
    plainheadsepline,% also on scrplain
}
\clearpairofpagestyles % Erase all current configurations (updated from deprecated \clearscrheadfoot)
% inner part of the header [titel_page]{other_pages}
\ihead[\coursename, \courseyear]{\coursename, \courseyear}
% outer part of the header [titel_page]{other_pages}
\ohead{\assignmentname}
% center part of the footer [titel_page]{other_pages}
\cfoot[\textup{Page \pagemark\ of \pageref{LastPage}}]{\textup{Page \pagemark\ of \pageref{LastPage}}}
% ------------- End of the headers and footers -------------- %



% ---------------- Setting up the title page ---------------- %
\title{\reportname}
\subtitle{An Evolutionary Algorithm for the Traveling Salesman Problem}
\author{\docauthor}
\date{\today}
% ------------------ End of the title page ------------------ %



% --------------- Packages and custom commands -------------- %
% Specify the input encoding of the source file as UTF8 
\usepackage[utf8]{inputenc}
% The T1 font encoding is an 8-bit encoding and fully supports words containing accented characters
\usepackage[T1]{fontenc}
% Load Latin modern fonts in outline format
\usepackage{lmodern}
% Load better typewriter font
\usepackage{inconsolata}
% For better hyphenation patterns
\usepackage[english]{babel}
% For improved micro-typography
\usepackage{microtype}
% Switch off extra space after a sentence
\frenchspacing
% To be able to iclude pictures
\usepackage{graphicx}
\usepackage{zref-totpages} % To get the total number of pages
\usepackage{lastpage} % For LastPage reference
\usepackage{mathtools,amsmath} % for advanced math formulas
 % To highlight source code
\usepackage{minted}
% Default options for all minted environments
\setminted{
           style=default, % for minted envronment
           autogobble, % automatically remove all common leading whitespace
           xleftmargin=10pt, % indentation to add (on the left) before the listing
           numbersep=6pt, % gap between numbers and start of line
           linenos, % enables line numbers
           mathescape % enables the usual math mode inside comments
}
\setmintedinline[c]{style=default} % for minted inlines
% Flexible handling of verbatim text
\usepackage{fancyvrb}
% For well-spaced lines and guidelines in tables
\usepackage{booktabs}
% To typeset algorithms or pseudocode in LaTeX 
\usepackage{algorithm}
\usepackage{algpseudocode}
% For plots and graphs
\usepackage{pgfplots}
\pgfplotsset{width=10cm,compat=1.18}

% To generate placeholder text - can be deleted when you have written real text
\usepackage{lipsum}
% ----------- End of packages and custom commands ----------- %


\begin{document}

\maketitle

% % It can be commented out if you do not want the table of contents
% \tableofcontents 


\section{Introduction}
\label{sec:intro}

\textcolor{red}{The introduction should briefly introduce the assignment and its purpose.}

\lipsum[4]


\section{The Traveling Salesman Problem}
\label{sec:problem_description}

\textcolor{red}{The second section should present the problem you tackle using your evolutionary approach. Overall, this section should include:}

{\color{red}
\begin{itemize}
    \item The mathematical formulation considered in your study. Some problems have a clear mathematical model (e.g., Travelling Salesman Problem), while others do not (e.g., $n$-Queens). Based on the problem you chose, search the literature and find a proper way to present the problem.
    \item One paragraph that briefly presents at least 3 published academic works where any evolutionary approach is used to solve the problem. It would be wise to cite here works that influenced your algorithm. This practice saves you time from looking for additional academic resources. You can find more information about reading and searching in the literature in \cite{zobel2014reading}.
    \item The motivation behind the evolutionary approach you decided to develop. A good practice would be to align the motivation with some literature gap found in the academic works you presented above. However, this is not mandatory. You can motivate your selection on the characteristics of the algorithm making it proper for the problem.
\end{itemize}
}

\lipsum[2]


\section{Algorithm}
\label{sec:algorithm}

\textcolor{red}{The third section should present the evolutionary approach you developed. You can divide this section into subsection. In any case, you should mention the following details:}

\textcolor{red}{\textbf{Evolutionary approach.} Clearly describe the algorithm you developed. You should clearly explain the evolutionary operators you used and what modifications you did to match the problem. It is extremely important to present also a pseudocode of your algorithm. An example is given in \ref{alg:pseudocode_example}, below. For more insight into presentation of algorithms, you can advise \cite{zobel2014algorithms}.}

{
\color{red}To typeset pseudocode in \LaTeX\ you can use one of the following options:
\begin{itemize}
    \item Choose ONE of the (\texttt{algpseudocode} OR \texttt{algcompatible} OR \texttt{algorithmic}) packages to typeset algorithm bodies, and the algorithm package for captioning the algorithm.
    \item The \texttt{algorithm2e} package.
\end{itemize}
You can find more information here: \url{https://www.overleaf.com/learn/latex/Algorithms}
}

\begin{algorithm}
\caption{Example of an algorithm's pseudocode}\label{alg:pseudocode_example}
\begin{algorithmic}
\Require $n \geq 0$
\Ensure $y = x^n$
\State $y \gets 1$
\State $X \gets x$
\State $N \gets n$
\While{$N \neq 0$}
\If{$N$ is even}
    \State $X \gets X \times X$
    \State $N \gets \frac{N}{2}$  \Comment{This is a comment}
\ElsIf{$N$ is odd}
    \State $y \gets y \times X$
    \State $N \gets N - 1$
\EndIf
\EndWhile
\end{algorithmic}
\end{algorithm}

\textcolor{red}{\textbf{Solution representation.} Clearly describe the solution representation you used. You can use figures to improve the comprehensibility of this part.}

\textcolor{red}{\textbf{Fitness function.} It is also very important to mention the fitness function you used. In many cases, the objective function of the problem is not the same as the fitness function used in an evolutionary algorithm. An example, following the principles of \cite{zobel2014mathematics}, is given below.}

\begin{equation}
    F = \sum_{i=1}^d x_i^2 
\end{equation}
where $x_i$ is the $i$-th gene (i.e., decision variable) in the solution and $d$ corresponds to the number of decision variables in the problem.

\textcolor{red}{\textbf{Note:} Change the section's title to match the name of the algorithm you developed for your assignment.}

\lipsum[3]


\section{Experimental part}
\label{sec:experimentation}

This section describes the setup of experiments conducted to evaluate and compare the performance of the greedy best-first search baseline with two evolutionary algorithm variants employing different crossover operators.

\subsection{Hardware and Software Specifications}

\subsubsection{Hardware Configuration}
The experiments were conducted on the following hardware configuration:
\begin{itemize}
    \item Processor: AMD Ryzen 7 7840HS (16 cores) @ 3.8GHz
    \item Memory: 48 GB RAM
    \item Operating System: Linux Ubuntu 24.04
\end{itemize}

\subsubsection{Software Environment}
The algorithms were implemented using the following software tools:
\begin{itemize}
    \item Programming Language: Python 3.12
    \item Libraries and Frameworks: 
    \begin{itemize}
        \item NumPy (numerical computations and array operations)
        \item Matplotlib (data visualization and convergence plots)
    \end{itemize}
    \item Development Environment: Visual Studio Code
\end{itemize}

\subsection{Experimental Procedure}

\subsubsection{Algorithm Configuration}

The evolutionary algorithm employs the following systematic design decisions:

\begin{itemize}
  \item{Selection Method}: Elite selection (fixed design choice ensuring the best individuals are preserved across generations)
  \item{Configurable Parameters}: The following parameters were varied during experimentation:
  \begin{itemize}
    \item{Population Size}
    \item{Stopping Criteria}
    \item{Crossover Rate}
    \item{Elite Rate}
  \end{itemize}
\end{itemize}

\subsubsection{Execution Protocol}
The experiments executed the algorithm following this protocol:
\begin{enumerate}
    \item \textbf{Initialization}: For evolutionary algorithms, the initial population was generated using random permutations of city sequences.  
    \item \textbf{Independent Runs}: Each algorithm configuration was run at least 25 times with different random seeds to account for stochastic variation and ensure statistical validity.
    \item \textbf{Termination Criteria}: The algorithm employ multiple configurable stopping criteria:
    \begin{itemize}
        \item \textbf{Iteration-based}: Maximum number of generations/iterations
        \item \textbf{Time-based}: Maximum execution time limit
        \item \textbf{Improvement-based}: Termination when improvement becomes less than a specified threshold
    \end{itemize}
    The specific stopping criteria (or combination of criteria) has been fine tuned during the parameter search process.
    \item \textbf{Results exporting}: The results were exported in a CSV file for each parameter set.
\end{enumerate}

\subsection{Evaluation Data}

\subsubsection{Problem Instances}
The algorithm were evaluated on TSP instances of varying sizes to assess scalability and performance:
\begin{itemize}
    \item \textbf{Problem Sizes}: 76, 136, and 280 cities (EIL76, PR136, A280 datasets)
    \item \textbf{Number of Instances}: 3 benchmark instances from TSPLIB
\end{itemize}

\subsubsection{Data Generation}
The experimental evaluation employs benchmark datasets from TSPLIB:

\textbf{Benchmark Datasets}: Standard TSP instances (EIL76, PR136, A280) from TSPLIB. These benchmark problems are parsed and stored as arrays of tuples containing city coordinates $(x, y)$.

This standardized approach allows for direct comparison with existing literature and ensures reproducibility of results across different implementations.

\subsubsection{Distance Metric}
The distance metric used for calculating distances between cities is configurable as a parameter:

\textbf{Euclidean Distance}: 
\begin{equation}
d(i,j) = \sqrt{(x_i - x_j)^2 + (y_i - y_j)^2}
\end{equation}

\textbf{Manhattan Distance}: 
\begin{equation}
d(i,j) = |x_i - x_j| + |y_i - y_j|
\end{equation}

where $(x_i, y_i)$ and $(x_j, y_j)$ represent the coordinates of cities $i$ and $j$ respectively.

\subsubsection{Performance Metrics}
The following metrics were used to evaluate and compare algorithm performance:

\begin{enumerate}
    \item \textbf{Best Cost}: The minimum tour length found across all runs for each algorithm configuration
    
    \item \textbf{Average Cost}: The mean tour length across all independent runs, providing insight into algorithm consistency
    
    \item \textbf{Standard Deviation}: The variability of solution quality across runs, indicating algorithm robustness
    
    \item \textbf{Average Execution Time}: The mean computational time required for each algorithm to complete
        
    \item \textbf{Baseline Gap (Best)}: The percentage difference between the best solution found and the baseline algorithm's best solution, calculated as:
    \begin{equation}
    \frac{EA_{best} - Baseline_{best}}{Baseline_{best}} \times 100\%
    \end{equation}
    
    \item \textbf{Baseline Gap (Average)}: The percentage difference between the average solution quality and the baseline algorithm's average solution quality, calculated as:
    \begin{equation}
    \frac{EA_{avg} - Baseline_{avg}}{Baseline_{avg}} \times 100\%
    \end{equation}
\end{enumerate}


\section{Results and Analysis}
\label{sec:results-analysis}

\subsection[Experiment 1]{Tuning population size / generation number ratio over cost}

This experiment was conducted to observe the effect of the population size / generation number ratio tuning on the cost. The results are shown in the following plots.

\subsubsection{Results for EIL76 dataset:}

\begin{tikzpicture}
\begin{axis}[
  width=0.95\textwidth, height=9cm,
  xlabel={Generation Nr / Population Size},
  ylabel={Cost},
  x tick label style={rotate=60,anchor=east},
  legend pos=north west,
  enlarge x limits=0.02,
  scaled x ticks=false,
  x tick label style={/pgf/number format/fixed, /pgf/number format/precision=3},
]

% average point (red circle) + whisker down to min
\addplot+[
  only marks,
  mark=o,
  color=red,
  thick,
  error bars/.cd,
    y dir=minus,   % draw whisker downward only
    y explicit,    % use given expression for its size
]
table[
  x={gen_pop_ratio},
  y=average_cost,
  y error expr={\thisrow{average_cost}-\thisrow{best_cost}},  % whisker length
  col sep=comma
]{tuning_ratio_over_cost_csv_files/eil76.csv};
\addlegendentry{average}

% bottom markers (min)
\addplot+[
  only marks,
  mark=square*,
  color=red,
  thick,
]
table[
  x={gen_pop_ratio},
  y=best_cost,
  col sep=comma
]{tuning_ratio_over_cost_csv_files/eil76.csv};
\addlegendentry{min}

\end{axis}
\end{tikzpicture}

- there is no noticeable evolution trend for average cost, values seem to fluctuate between 1000 and 1200
- there are several variations between average cost and best cost, values variating from a difference of 100 to even 400 at some point

\subsubsection{Results for PR136 dataset:}

\begin{tikzpicture}
\begin{axis}[
  width=0.95\textwidth, height=9cm,
  xlabel={Generation Nr / Population Size},
  ylabel={Cost},
  x tick label style={rotate=60,anchor=east},
  legend pos=north west,
  enlarge x limits=0.02,
  scaled x ticks=false,
  x tick label style={/pgf/number format/fixed, /pgf/number format/precision=3},
]

% average point (red circle) + whisker down to min
\addplot+[
  only marks,
  mark=o,
  color=red,
  thick,
  error bars/.cd,
    y dir=minus,   % draw whisker downward only
    y explicit,    % use given expression for its size
]
table[
  x={gen_pop_ratio},
  y=average_cost,
  y error expr={\thisrow{average_cost}-\thisrow{best_cost}},  % whisker length
  col sep=comma
]{tuning_ratio_over_cost_csv_files/pr136.csv};
\addlegendentry{average}

% bottom markers (min)
\addplot+[
  only marks,
  mark=square*,
  color=red,
  thick,
]
table[
  x={gen_pop_ratio},
  y=best_cost,
  col sep=comma
]{tuning_ratio_over_cost_csv_files/pr136.csv};
\addlegendentry{min}

\end{axis}
\end{tikzpicture}

- slight trend of the cost going upwards, leaving aside some exceptional cases
- the difference between average and best cost is pretty uniform for the majority of cases, variating between 0.5*10(pow)5 and 1*10(pow)5 
- to give context, the majority of average costs for each point vary around 3.5*10(pow)5 and 4.5*10(pow)5

\subsubsection{Results for A280 dataset:}

\begin{tikzpicture}
\begin{axis}[
  width=0.95\textwidth, height=9cm,
  xlabel={Generation Nr / Population Size},
  ylabel={Cost},
  x tick label style={rotate=60,anchor=east},
  legend pos=north west,
  enlarge x limits=0.02,
  scaled x ticks=false,
  x tick label style={/pgf/number format/fixed, /pgf/number format/precision=3},
]

% average point (red circle) + whisker down to min
\addplot+[
  only marks,
  mark=o,
  color=red,
  thick,
  error bars/.cd,
    y dir=minus,   % draw whisker downward only
    y explicit,    % use given expression for its size
]
table[
  x={gen_pop_ratio},
  y=average_cost,
  y error expr={\thisrow{average_cost}-\thisrow{best_cost}},  % whisker length
  col sep=comma
]{tuning_ratio_over_cost_csv_files/a280.csv};
\addlegendentry{average}

% bottom markers (min)
\addplot+[
  only marks,
  mark=square*,
  color=red,
  thick,
]
table[
  x={gen_pop_ratio},
  y=best_cost,
  col sep=comma
]{tuning_ratio_over_cost_csv_files/a280.csv};
\addlegendentry{min}

\end{axis}
\end{tikzpicture}

- there seems to be a general trend of the cost going upwards
- differences between average cost and best cost seem to be approximately the same for almost every point, few exceptions aside

\subsubsection{Experiment results analysis}

During this experiment, we observed that the cost decreases as the population size / generation number ratio increases. This is expected because a larger population size allows for more exploration of the search space, which leads to better solutions.

Conclusion: 
- we do not obtain a specific trend from this ratio based on these three datasets, seems either fluctuating or going upwards
- we know that usually a higher number of generations hands a lower cost up to some point decided somehow by the population size and the number of nodes, then it stalls to a certain value
- the test seems to be inconclusive

\section{Conclusions}
\label{sec:conclusions}

\textcolor{red}{In this section you should provide a concise summary of what has been done, the obtained results and some recommendations on how this study could be extended.}

\lipsum[4]

\bibliographystyle{ieeetr}
\bibliography{references}

\end{document}
